\section{More Background on Privacy Amplification and Correlated Noise}\label{app:background}

\mypar{Privacy amplification} Though randomness in data processing reduces the variance of noise required, the main challenge with privacy amplification is giving tight privacy analysis of the amplified mechanism, which for DP-SGD usually requires computing a divergence between two mixtures of Gaussians. There are several forms of privacy amplification, but the two most popular are sampling \cite{DP-DL, BBG18} and shuffling \cite{balle2018improving, erlingsson2019private, PAS, feldman22hiding}.

Privacy amplification by (Poisson) sampling assumes the batches are formed by including each element independently with some probability $p$. Because the sampling randomness is independent across rounds, tight privacy analyses of DP training with independent noise (DP-SGD) are well-known \cite{koskela21tight}, and supported by open-source libraries  \cite{dp_accounting}. However, Poisson sampling is often impractical to implement, especially when working with larger datasets~\cite{ponomareva23dpfy}.

Privacy amplification by shuffling in the context of DP-SGD usually assumes the batches are formed by shuffling the dataset (a single time) and then forming fixed-size batches using the shuffled order. In some limited settings its privacy analysis is well-understood \cite{feldman22hiding}, but in the setting of DP-SGD, for example, tight privacy analyses remain elusive. Technically, the difficulty is that the participations of different examples are dependent, unlike with Poisson sampling. In addition to being challenging to analyze, \citet{chua2024private} demonstrated that shuffling often gives much smaller improvements in privacy compared to sampling. Despite these difficulties, as \citet{chua2024private} note, shuffling is used much more widely in practice since it can be implemented, e.g., in a single pre-processing pass over the dataset, whereas sampling requires random access to potentially very large datasets which is often infeasible.

\mypar{Correlated noise} 
\citet{kairouz2021practical} proposed DP-FTRL, the first DP training algorithm with correlated noise, where the update at each step are the rows of $\bfx + \bfz_{tree}$ where $\bfz_{tree}$ is sampled using the binary tree mechanism of \cite{Dwork-continual}.~\citet{denisov2022improved} showed this was a special case of the updates being rows of $\bfx + \bfC^{-1} \bfz$, giving a more expressive definition for $\bfC^{-1}$ called the ``correlation matrix''. They showed these matrices could be optimized via a convex optimization program to maximize noise cancellation (i.e., minimize total noise variance).~\citet{choquette2022multi} proposed a ``multi-participation'' generalization, leading to the first correlated noise DP training algorithm to (substantially) outperform DP-SGD with practical $\epsilon$'s.~\citet{choquette2024amplified} showed that correlated noise algorithms Pareto-dominated DP-SGD in privacy-utility tradeoffs;~\citet{choquette2023correlated} showed that these algorithms provably outperformed DP-SGD.

For simplicity, these works assume $\bfC$ is non-negative and lower-triangular. The privacy guarantee of correlated noise is the same as that of the matrix mechanism $\bfC \bfx + \bfz$ by post-processing, hence based on how batches are formed we can usually give privacy guarantees in terms of an appropriate norm of $\bfC$. To optimize the matrix $\bfC$, the convex program usually minimizes the root-mean-squared-error (RMSE) of the prefix sums of $\bfx$, given by $\ltwo{\bfA \bfC^{-1}} \cdot \sigma$ where $\bfA$ is the lower-triangular all-ones matrix and $\sigma$ is the noise multiplier needed for $\bfC \bfx + \bfz$ to satisfy the target privacy guarantee.